%% kapa: exam=1 round=30 year=2019 period=1 form=A subject=civil_law subject_ko=민법 seq=27 number=27
%% source: Q-Net 감정평가사 기출문제 (cst00309) artlSeq=5205205 file=extracted/1008/5205205/2019년 제30회 감정평가사 1차시험 문제지/2019년 제30회 감정평가사 1차시험 문제지/2019년도 제30회 감정평가사 1차 1교시 A형.hwp
%% conversion: kordoc 4.24.51 -> tools/gen1.py (auto); review=proofread-sonnet; flags=none
\begin{question}{27}
물권의 소멸에 관한 설명으로 옳지 않은 것은? (다툼이 있으면 판례에 따름)
\begin{choices}
  \item 물건이 멸실되더라도 물건의 가치적 변형물이 남아 있는 경우에는 담보물권은 그 가치적 변형물에 미친다.
  \item 지역권은 소멸시효의 대상이 될 수 있다.
  \item 부동산에 대한 합유지분의 포기는 형성권의 행사이므로 등기하지 않더라도 포기의 효력이 생긴다.
  \item 점유권과 본권이 동일인에게 귀속하더라도 점유권은 소멸하지 않는다.
  \item 근저당권자가 그 저당물의 소유권을 취득하면 그 근저당권은 원칙적으로 혼동에 의하여 소멸하지만, 그 뒤 그 소유권 취득이 무효인 것이 밝혀지면 소멸하였던 근저당권은 당연히 부활한다.
\end{choices}
\end{question}
\answer{3}
